\subsection{Sprint 2}

Antes mesmo do início da \textit{sprint}, eu e os demais \ac{ages} IV do time revisamos todas as tarefas que já haviam sido criadas, de modo que ficassem alinhadas tanto com as regras de negócio quanto com o modelo de banco de dados, facilitando o desenvolvimento. Dividimos entre duas e três tarefas para cada um revisar, conforme a complexidade. Enquanto revisava as que ficaram sob minha responsabilidade, referentes à execução de simulados, percebi no código um detalhe do \textit{frontend} que poderia impactar a \textit{sprint}: alguns componentes que deveriam ter sido criados na \textit{sprint} anterior, para serem reutilizados nesta, haviam sido implementados de forma \textit{hard coded}, e não como componentes reaproveitáveis. Avisei os demais \ac{ages} IV do time e ajustei a tarefa desta \textit{sprint} para que isso fosse corrigido, seguindo as boas práticas de engenharia de software.

Nesta \textit{sprint}, definimos que as entregas seriam a outra \textit{feature} central da aplicação, a parte de criação e execução de simulados, além do \textit{onboarding} do usuário e do bloqueio de pagamentos (\textit{paywall}). Para a execução de simulados, optamos por reutilizar parte do código já desenvolvido para as questões, já que, na prática, ambos cumprem o mesmo propósito. A única diferença é que, nas questões de treino, o usuário recebe o \textit{feedback} de forma imediata, enquanto no simulado isso ocorre apenas ao finalizá-lo.

Decidimos também alterar um pouco a metodologia de trabalho. Em vez de distribuir tarefas específicas para cada \textit{squad}, passamos a distribuí-las de forma individual, deixando a \textit{squad} apenas como um norte, indicando a quem a pessoa deveria recorrer primeiro caso precisasse de ajuda. Assim, no meu time, as tarefas ficaram divididas entre os épicos: pegamos atividades de execução de simulados, de \textit{paywall} e algumas telas do \textit{dashboard} de administração.

No início da \textit{sprint}, tivemos um problema de comunicação no time. A forma como a tabela de simulados seria utilizada foi alterada, mas a mudança demorou bastante para ser \textit{mergeada} no repositório, o que gerou retrabalho, pois as demais tarefas já estavam em desenvolvimento quando souberam dessa informação. Por conta disso, junto com os \ac{ages} III e o pessoal que estava desenvolvendo as tarefas de simulados, reformulamos alguns pontos da regra de negócio: removemos o estado de ``pausado'' e passamos a permitir apenas um simulado ativo por \textit{id} de simulado. Ou seja, se um usuário possui um simulado ativo e deseja refazê-lo, o \textit{backend} entende que ele está encerrando a tentativa em aberto e iniciando uma nova.

Ainda na primeira semana de desenvolvimento, junto com os \ac{ages} III responsáveis pela \ac{ui}/\ac{ux} do projeto, marcamos uma reunião com as \textit{stakeholders} para apresentar o que havíamos imaginado para as telas que elas apontaram precisar de ajustes: a tela inicial de treino, em que o usuário escolhe a disciplina, e a tela de trilha, em que ele escolhe o tópico. Elas aprovaram as alterações, mas, como o escopo desta \textit{sprint} já estava definido, optamos por não as incorporar agora, deixando-as para a \textit{Sprint} 3.

Entre o meio e o fim da \textit{sprint}, precisei me ausentar das aulas por conta de uma cirurgia. Ainda assim, procurei manter contato tanto com o meu time quanto com os demais \ac{ages} IV para acompanhar o andamento do projeto. Quando retornei, estávamos bastante atrasados para a entrega, já nos últimos dias de \textit{sprint}. Em tese, o \textit{code freeze} já deveria ter ocorrido, mas, se o respeitássemos, não teríamos nada a entregar. Por isso, decidimos que cada um enviaria o que conseguisse finalizar e os \ac{ages} III iriam corrigindo conforme fosse possível.

Na terça-feira, véspera da apresentação, me reuni com os demais \ac{ages} IV para preparar a apresentação, e foi quando notamos que quase todas as tarefas entregues estavam com \textit{bugs}. Como ainda estava afastado do trabalho e estava com bastante tempo disponivel para me dedicar ao projeto naquele período, prontifiquei-me a corrigir o que fosse possível até a hora da apresentação. Durante todo o dia da apresentação para as \textit{stakeholders}, fiz ajustes no código para que conseguíssemos entregar algo a elas. No fim, consegui corrigir boa parte dos \textit{bugs} e, mais uma vez, as \textit{stakeholders} ficaram bastante satisfeitas com a entrega.

Na retrospectiva, surgiram vários pontos sobre o \textit{code freeze} não ter sido respeitado, sobre a demora na correção dos \ac{pr} por parte dos \ac{ages} III e sobre a alteração no banco de dados, que demorou a ser \textit{mergeada} na \textit{develop} e não foi devidamente comunicada aos demais colegas. Também foi trazido que os \ac{ages} III estavam se sentindo sobrecarregados com a quantidade de tarefas e revisões que precisavam fazer. Com base nisso, junto com a equipe, montamos alguns pontos de ação para evitar que esses problemas voltassem a ocorrer. Estabelecemos que, sempre que algo fosse alterado no banco de dados de forma que pudesse influenciar o trabalho de outras pessoas, isso deveria ser notificado no Discord, no canal designado para o banco de dados. Também avisamos aos \ac{ages} III que iríamos rever como poderíamos reduzir essa sobrecarga deles nas próximas \textit{sprints}. Por fim, reforçamos que todo o código de \textit{frontend} deveria ser testado utilizando o Expo, e não apenas no navegador, pois assim, durante o desenvolvimento, conseguiríamos ter uma ideia melhor de como ele ficaria no celular.
